\documentclass[DIV=calc]{scrartcl}
\usepackage[utf8]{inputenc}
\usepackage[T1]{fontenc}
\usepackage[ngerman]{babel}
\usepackage{graphicx}
\usepackage[draft, markup=underlined]{changes}
\usepackage{csquotes}
\usepackage{eurosym}

\usepackage{ulem}
%\usepackage[dvipsnames]{xcolor}
\usepackage{paralist}
%\usepackage{fixltx2e}
%\usepackage{ellipsis}
\usepackage[tracking=true]{microtype}

\usepackage{lmodern}              % Ersatz fuer Computer Modern-Schriften
%\usepackage{hfoldsty}

%\usepackage{fourier}             % Schriftart
\usepackage[scaled=0.81]{helvet}     % Schriftart

\usepackage{url}
%\usepackage{tocloft}             % Paket für Table of Contents
\def\UrlBreaks{\do\a\do\b\do\c\do\d\do\e\do\f\do\g\do\h\do\i\do\j\do\k\do\l%
\do\m\do\n\do\o\do\p\do\q\do\r\do\s\do\t\do\u\do\v\do\w\do\x\do\y\do\z\do\0%
\do\1\do\2\do\3\do\4\do\5\do\6\do\7\do\8\do\9\do\-}%

\usepackage{xcolor}
\definecolor{urlred}{HTML}{660000}

\usepackage{hyperref}
\hypersetup{colorlinks=false}

%\usepackage{mdwlist}     % Änderung der Zeilenabstände bei itemize und enumerate
% \usepackage{draftwatermark} % Wasserzeichen ``Entwurf''
% \SetWatermarkText{Antrag}

\parindent 0pt                 % Absatzeinrücken verhindern
\parskip 12pt                 % Absätze durch Lücke trennen

\setlength{\textheight}{23cm}
\usepackage{fancyhdr}
\pagestyle{fancy}
\fancyhead{} % clear all header fields
\cfoot{}
\lfoot{Zusammenkunft aller Physik-Fachschaften}
\rfoot{www.zapfev.de\\stapf@zapf.in}
\renewcommand{\headrulewidth}{0pt}
\renewcommand{\footrulewidth}{0.1pt}
\newcommand{\gen}{*innen}
\addto{\captionsngerman}{\renewcommand{\refname}{Quellen}}

%%%% Mit-TeXen Kommandoset
\usepackage[normalem]{ulem}
\usepackage{xcolor}
\usepackage{xspace} 

\newcommand{\replace}[2]{
    \sout{\textcolor{blue}{#1}}~\textcolor{blue}{#2}}
\newcommand{\delete}[1]{
    \sout{\textcolor{red}{#1}}}
\newcommand{\add}[1]{
    \textcolor{blue}{#1}}

\newif\ifcomments
\commentsfalse
%\commentstrue

\newcommand{\red}[1]{{\ifcomments\color{red} {#1}\else{#1}\fi}\xspace}
\newcommand{\blue}[1]{{\ifcomments\color{blue} {#1}\else{#1}\fi}\xspace}
\newcommand{\green}[1]{{\ifcomments\color{green} {#1}\else{#1}\fi}\xspace}

\newcommand{\repl}[2]{{\ifcomments{\color{red} \sout{#1}}{\color{blue} {\xspace #2}}\else{#2}\fi}}
%\newcommand{\repl}[2]{{\color{red} \sout{#1}\xspace{\color{blue} {#2}}\else{#2}\fi}\xspace}

\newcommand{\initcomment}[2]{%
	\expandafter\newcommand\csname#1\endcsname{%
		\def\thiscommentname{#1}%
		\definecolor{col}{rgb}{#2}%
		\def\thiscommentcolor{col}%
}}

% initcomment Name RGB-color
\initcomment{Philipp}{0, 0.5, 0}

%\renewcommand{\comment}[1]{{\ifcomments{\color{red} {#1}}{}\fi}\xspace}

\renewcommand{\comment}[2][\nobody]{
	\ifdefined#1
	{\ifcomments{#1 \expandafter\color{\thiscommentcolor}{\thiscommentname: #2}}{}\fi}\xspace
	\else
	{\ifcomments{\color{red} {#2}}{}\fi}\xspace
	\fi
}

\newcommand{\zapf}{ZaPF\xspace}

\let\oldgrqq=\grqq
\def\grqq{\oldgrqq\xspace}

\setlength{\parskip}{.6em}
\setlength{\parindent}{0mm}

%\usepackage{geometry}
%\geometry{left=2.5cm, right=2.5cm, top=2.5cm, bottom=3.5cm}

% \renewcommand{\familydefault}{\sfdefault}




\begin{document}

\hspace{0.87\textwidth}
\begin{minipage}{120pt}
	\vspace{-1.8cm}
	\includegraphics[width=80pt]{logo.pdf}
	\centering
	\small Zusammenkunft aller Physik-Fachschaften
\end{minipage}

\begin{center}
  \huge{Studentischer Tarifvertrag, jetzt unterstützen!}\vspace{.25\baselineskip}\\
  \normalsize
\end{center}
\vspace{1cm}

%%%% Metadaten %%%%

% \paragraph{Adressierte:} ASten und alle (Physik-)Fachschaften, LAK (Landes-ASten-Konferenz), MeTaFa, fzs, TVStud, GEW-Studis, Verdi, FAU



% \paragraph{Antragstellende:} Chris (Uni Marburg), Felix (TU Darmstadt), Jonathan (Uni Bremen), Fritz und Leo (FUB)

% %%%% Text des Antrages zur veröffentlichung %%%%

% \section*{Antragstext}

Die Studie “Jung, Akademisch, Prekär” (Hoffmann, Hopp et Al., 2022)\footnote{Hoffmann, Hopp et Al., 2022: \url{https://www.iaw.uni-bremen.de/f/b1d9874527.pdf}} zeigt: Die Situation studentisch Beschäftigter an Hochschulen ist prekär. Drei Viertel der studentisch Beschäftigten sind von Armut gefährdet [S.69, Abb.31], mit einer durchschnittlichen Vertragslaufzeit von unter 6 Monaten [S.70ff, Abb.32] werden Studierende unnötig unter Stress gesetzt. Unnötig, da sie sich oft ohnehin in einer Kettenbefristung befinden [S.76, Abb.36]. Jedes Semester ein neuer Vertrag - auf der selben Stelle und nie können sich die Beschäftigten sicher sein, ob sie sich in kurzfristiger Zukunft noch ihre eigene Ausbildung leisten können. Hinzu kommen fehlende Informationen über die eigenen Arbeitnehmer:innenrechte [S.92] oder eine demokratische Partizipationsrechte bei der aktiven und passiven Wahl von Personalvertretungen [S.33, Abb.6].

Dabei ist die Lage jetzt günstig: In den Bundesländern steht es 10:6, eine Mehrheit der Länder ist für einen TVStud - einen Tarifvertrag für Studierende, worüber ab Herbst verhandelt wird. Für eine Verbesserung der Situation studentisch Beschäftigter ist ein Tarifvertrag essentiell, sie müssen sich vernetzen\footnote{Kontakte der Bundes- und Lokalinitiativen: \url{https://tvstud.de/mitmachen/netzwerk/}}. Dadurch können eine angemessene Bezahlung, Vertretung studentischer Interessen in Personalräten und das Problem von Kettenbefristungen effektiver bekämpft werden.

Zum Beispiel müssen Arbeitgeber:innen den Studierenden dann Verträge mit einer Mindestvertragslaufzeit anbieten und es liegt im Ermessen der Beschäftigten, diese Mindestlaufzeit in Anspruch zu nehmen oder situationsbedingt auch eine kürzere Laufzeit zu vereinbaren. Damit ist es für Studierende nur von Vorteil, da eine langfristigere finanzielle Sicherheit gegeben ist und sie mehr Entscheidungsspielraum haben.

Die ZaPf fordert daher alle (Physik-)Studierenden, Fachschaftsvertretungen und ASten auf, in Zusammenarbeit mit ihren lokalen TVStud-Initiativen über die Lage studentisch Beschäftigter zu informieren sowie die Initiative selbst zu unterstützen. Dabei sollen Treffen lokaler Initiativen, das Schaffen von Anlaufstellen bei den ASten, der bundesweite Austausch und die Möglichkeit, TVStud-Botschafter:in zu sein, beworben werden. Hierzu gibt die gemeinsame Resolution verschiedener MeTaFa-Fachschaften "Forderung: Tarifverträge für studentische Beschäftigte"\footnote{\href{https://www.fzs.de/buendnisse-zusammenarbeit-und-mitgliedschaften/meta-tagung-der-fachschaften-metafa/forderung-tarifvertrag-fuer-studentische-beschaeftigte-tvstud-2/}{Zu finden auf der Webseite des fzs, abgerufen am 18.05.2023}} eine Basis.

Damit die Forderungen studentischer Beschäftigter bis zur Tarifverhandlung gehört werden, müssen wir jetzt aktiv werden!



% \section*{Begründung}
% Erfolgt mündlich.

\vspace{1cm} 

\vfill
\begin{flushright}
	Verabschiedet am 01. Mai 2023 \\
	auf der ZaPF in Berlin.
\end{flushright}

\end{document}
